\section{They Shall See His Face}

The Church ends her year with her eyes lifted. On the last Sunday after Pentecost the 1962 Missal reads the discourse of the Lord on the end of the age, with its promise that the nations ``shall see the Son of man coming in the clouds of heaven with much power and majesty''; and on the First Sunday of Advent, beginning the year again, it reads the same promise in Luke with a command attached: ``look up and lift up your heads, because your redemption is at hand.''\endnote{Matt 24:30; Luke 21:28. Douay--Rheims (Challoner). \emph{Missale Romanum} (1962), Dominica XXIV et ultima post Pentecosten, Evangelium (Matt 24:15--35); Dominica I Adventus, Evangelium (Luke 21:25--33).} The last weeks of the year are full of judgment. The season, though, opened with All Saints, when both Missals read the Beatitudes, among them the promise of the vision: ``Blessed are the clean of heart: for they shall see God.'' The 1962 Missal sings that verse again at the Communion of the day, while the faithful receive Christ; the Postconciliar Missal sets beside the Gospel a reading from the first epistle of John.\endnote{Matt 5:8. \emph{Missale Romanum} (1962), In festo Omnium Sanctorum, Evangelium (Matt 5:1--12) and Communio (Matt 5:8--10). \emph{Ordo lectionum Missae}, editio typica altera (1981), n.~667: Apoc 7:2--4, 9--14; 1 John 3:1--3; Matt 5:1--12a.}

In it John says what Christians are now, and why they cannot yet know what they will be:

\begin{quote}
Behold what manner of charity the Father hath bestowed upon us, that we should be called and should be the sons of God. \ldots\ Dearly beloved, we are now the sons of God: and it hath not yet appeared what we shall be. We know that when he shall appear we shall be like to him: because we shall see him as he is.\endnote{1 John 3:1--2, Douay--Rheims (Challoner).}
\end{quote}

John first admits what is not known: \emph{It hath not yet appeared what we shall be.} Then he gives a reason, and the reason is a promise: we shall be like him \emph{because} we shall see him as he is. This life is lived by faith and does not give that sight, so no one can yet describe the likeness it will bring. The apostle draws the practical consequence in the next verse: ``every one that hath this hope in him sanctifieth himself, as he also is holy.''\endnote{1 John 3:3, Douay--Rheims (Challoner).}

\subsection{The city described by what it lacks}

The last two chapters of Scripture describe the city of God at length, and nearly every line of the description takes something away. In the new heaven and the new earth ``the sea is now no more.'' ``God shall wipe away all tears from their eyes: and death shall be no more. Nor mourning, nor crying, nor sorrow shall be any more, for the former things are passed away.'' ``I saw no temple therein. For the Lord God Almighty is the temple thereof, and the Lamb.'' ``The city hath no need of the sun, nor of the moon, to shine in it.'' ``The gates thereof shall not be shut by day: for there shall be no night there.'' ``There shall be no curse any more.''\endnote{Apoc 21:1, 4, 22, 23, 25; 22:3, Douay--Rheims (Challoner).} The Apocalypse speaks of heaven mainly by subtraction, as one must speak of something the hearer has no positive notion of: it removes, one by one, the conditions of this life that wound.

Among the negations stands one clause that takes nothing away:

\begin{quote}
And they shall see his face: and his name shall be on their foreheads.\endnote{Apoc 22:4, Douay--Rheims (Challoner).}
\end{quote}

Several of the subtractions are explained by that presence. There is no temple because the whole city stands before God, and there is no sun because ``the glory of God hath enlightened it: and the Lamb is the lamp thereof.''\endnote{Apoc 21:23, Douay--Rheims (Challoner).} The gates stand open because nothing is left to fear, and the name written on the foreheads marks the blessed as belonging to the one whose face they see.

Speaking by subtraction also keeps the reader from furnishing heaven with what he already knows. Alphonsus, setting the joys of paradise before his readers, gives them a homely warning about the measure of their imagination:

\begin{quote}
Were a horse capable of reasoning, he would, if he expected a rich feast from his master, imagine it to consist in excellent hay and oats; for these are the only species of food of which he has any idea. It is thus we form our notions of the happiness of Heaven.\endnote{Alphonsus Liguori, \emph{Preparation for Death}, Consideration XXIX, first point (Dublin: James Duffy, 1843).}
\end{quote}

Scripture's own images are another matter, because they are revealed, and the Catechism gathers them: ``Scripture speaks of it in images: life, light, peace, wedding feast, wine of the kingdom, the Father's house, the heavenly Jerusalem, paradise.''\endnote{\emph{Catechism of the Catholic Church} 1027.} A picture can show a city, a feast, a garden and a multitude in white. It cannot show God, who has no shape, though the city, the feast and the multitude are there for the sake of seeing him.

\subsection{Heaven is a Person, seen}

The Church rightly speaks of heaven as a place; it is also the reward of the just, and a reunion. The Catechism names what comes first: ``The beatific vision, in which God opens himself in an inexhaustible way to the elect, will be the ever-flowing well-spring of happiness, peace, and mutual communion.''\endnote{\emph{Catechism of the Catholic Church} 1045.}

Since the vision cannot be pictured, what the Church asks of her children at the end of her year is less to imagine heaven than to desire it, with the desire of the psalm: ``As the hart panteth after the fountains of water; so my soul panteth after thee, O God. My soul hath thirsted after the strong living God; when shall I come and appear before the face of God?''\endnote{Ps 41:2--3, Douay--Rheims (Challoner).} And when one of her children is carried out of the church to the grave, she sings, asking the angels to lead the soul into paradise: \latin{In paradisum deducant te Angeli}.
